\chapter{Mechanized Proof Checking}
\label{ch:mechanization}

\section{Kernel, elaboration, and automation}
\label{sec:mechanization-01}

\begin{definition}[Proof term]
Relative to a type-theoretic proof assistant, a \emph{proof term} is a term
whose type is the formal statement proved, according to the assistant's kernel
type theory.
\end{definition}

\begin{definition}[Kernel]
A \emph{kernel} is the component that checks declarations and proof terms
against the primitive typing and conversion rules of the implemented
foundation.
\end{definition}

\begin{definition}[Elaboration]
\emph{Elaboration} is the transformation of surface-language input containing
notation, omitted arguments, coercions, and tactic output into a kernel-level
term with all information required for checking.
\end{definition}

\begin{definition}[Tactic]
A \emph{tactic} is a program that transforms a proof state or constructs a
candidate proof term.  A tactic is not an additional inference rule when its
output is checked by the kernel.
\end{definition}

Lean 4, Rocq, and Isabelle expose different foundations and user languages.
Their official manuals distinguish surface syntax from kernel or framework
logic \cite{LeanManual2026,LeanTPIL,DeMouraUllrich2021,RocqManual2026,Isabelle2025}.
A theorem
checked in one assistant is relative to the definitions, axioms, imports, and
kernel of that development.  For a survey of the mechanization of
mathematics, see \cite{Avigad2018}.

